\subsection{Option Pricing Setup}\label{sec:pricing-setup}
% The generator approximates the conditional distribution of the next daily log
% return. Option pricing requires the distribution of the underlying price at
% maturity, which lies several weeks ahead. Paths were therefore generated
% recursively.

% The condition was initialized with the $l$ realized returns preceding the
% valuation date. At each step, one return was drawn from the generator and
% appended to the condition, while the oldest observation was removed. The step
% was repeated until the number of generated returns equalled the number of
% trading days to maturity. After $l$ steps the condition consists of generated
% returns only. Realized returns were not used after initialization, since they
% are not available at the valuation date.

% The simulated price at maturity follows from compounding the generated
% returns,
% \begin{equation}
%     \hat{S}_T = S_0 \exp\left( \sum_{k=1}^{T} \hat{r}_k \right),
% \end{equation}
% where $S_0$ is the settlement price at the valuation date and $T$ is the
% number of trading days to maturity. The procedure was repeated with
% independent noise samples to obtain a distribution of $\hat{S}_T$.

% The recursive scheme places the generator outside its training regime. The
% networks were trained on conditions consisting of realized returns, whereas
% after $l$ steps the condition consists of generated returns only. Deviations
% in the generated distribution therefore accumulate along the path. The
% consequences for the simulated terminal distribution are examined in
% Sec.~\ref{sec:pricing-results}. 

% \todo{Decide: valuation date and expiry for the pricing benchmark. Requirements
% --- front-month futures price exists (Bloomberg series ends 2026-04-24), enough
% listed strikes near the money, and relative spreads tight enough to trust
% (BBO note: usable mid prices begin in 2025). Candidate window: 2025 to April
% 2026, inside the test split. Query the master table for candidates.}